\documentclass[11pt]{article}

\usepackage[T1]{fontenc}
\usepackage[utf8]{inputenc}
\usepackage{lmodern}
\usepackage[margin=1.08in]{geometry}
\usepackage{amsmath,amssymb,mathtools,amsthm}
\usepackage{mathrsfs}
\usepackage{microtype}
\usepackage{enumitem}
\usepackage{booktabs}
\usepackage{array}
\usepackage{xcolor}
\usepackage{hyperref}
\usepackage[capitalize,noabbrev]{cleveref}

\hypersetup{
  colorlinks=true,
  linkcolor=blue!48!black,
  citecolor=green!35!black,
  urlcolor=blue!55!black,
  pdftitle={Relative Exponential Fixed Parts and Functional Periods: An Exponential Nori--Ayoub Theorem},
  pdfauthor={Christopher D. Long; Antoine-Auguste Le Blanc},
  pdfsubject={Exponential motives, perverse Nori motives, fixed parts, differential Galois groups, functional periods},
  pdfkeywords={exponential motive, Nori motive, Ayoub, fixed part, Fourier-Laplace transform, rapid decay, functional period}
}

\numberwithin{equation}{section}
\setlength{\emergencystretch}{3em}
\allowdisplaybreaks

\newtheorem{theorem}{Theorem}[section]
\newtheorem{proposition}[theorem]{Proposition}
\newtheorem{lemma}[theorem]{Lemma}
\newtheorem{corollary}[theorem]{Corollary}
\newtheorem{conjecture}[theorem]{Conjecture}
\newtheorem{question}[theorem]{Question}
\theoremstyle{definition}
\newtheorem{definition}[theorem]{Definition}
\newtheorem{hypothesis}[theorem]{Hypothesis}
\theoremstyle{remark}
\newtheorem{remark}[theorem]{Remark}
\newtheorem{warning}[theorem]{Warning}

\newcommand{\Aone}{\mathbb A^1}
\newcommand{\Gm}{\mathbb G_m}
\newcommand{\Ga}{\mathbb G_a}
\newcommand{\Q}{\mathbb Q}
\newcommand{\C}{\mathbb C}
\newcommand{\kbar}{\overline{k}}
\newcommand{\MNQ}{\mathsf{MN}_{\Q}}
\newcommand{\MN}{\mathsf{MN}}
\newcommand{\EMN}{\mathsf{EM}_{\mathrm N}}
\newcommand{\EMH}{\mathsf{EM}_{\mathrm H}}
\newcommand{\MHM}{\mathsf{MHM}}
\newcommand{\Hol}{\mathsf{Hol}}
\newcommand{\Conn}{\mathsf{Conn}}
\newcommand{\Loc}{\mathsf{Loc}}
\newcommand{\Perv}{\mathsf{Perv}}
\newcommand{\Tcal}{\mathcal T}
\newcommand{\Ccal}{\mathcal C}
\newcommand{\Pcal}{\mathcal P}
\newcommand{\Ecal}{\mathcal E}
\newcommand{\Fcal}{\mathcal F}
\newcommand{\Ocal}{\mathcal O}
\newcommand{\Rcal}{\mathcal R}
\newcommand{\Mcal}{\mathcal M}
\newcommand{\Hom}{\operatorname{Hom}}
\newcommand{\End}{\operatorname{End}}
\newcommand{\Aut}{\operatorname{Aut}}
\newcommand{\im}{\operatorname{im}}
\newcommand{\Spec}{\operatorname{Spec}}
\newcommand{\Pic}{\operatorname{Pic}}
\newcommand{\gr}{\operatorname{gr}}
\newcommand{\Fix}{\operatorname{Fix}}
\newcommand{\real}{\operatorname{real}}
\newcommand{\per}{\operatorname{per}}
\newcommand{\Isom}{\operatorname{Isom}}
\newcommand{\Gal}{\operatorname{Gal}}
\newcommand{\trdeg}{\operatorname{trdeg}}
\newcommand{\dR}{\mathrm{dR}}
\newcommand{\dRB}{\mathrm{dRB}}
\newcommand{\rd}{\mathrm{rd}}
\newcommand{\Betti}{\mathrm B}
\newcommand{\mot}{\mathrm{mot}}
\newcommand{\diff}{\mathrm{diff}}
\newcommand{\const}{\mathrm{const}}
\newcommand{\PV}{\mathrm{PV}}
\newcommand{\St}{\mathrm{St}}
\newcommand{\rank}{\operatorname{rank}}
\newcommand{\id}{\operatorname{id}}

\title{Relative Exponential Fixed Parts and Functional Periods:\\
An Exponential Nori--Ayoub Theorem}
\author{Christopher D. Long \and Antoine-Auguste Le Blanc\thanks{Le Blanc cryptographic author identifier: \texttt{b6048125\allowbreak{}30b3535b\allowbreak{}23d6dbb0\allowbreak{}f5abe32d\allowbreak{}1ed1a9de\allowbreak{}35fb4f7d\allowbreak{}04d62cff\allowbreak{}08ea19c8}}}
\date{August 2026}

\begin{document}
\maketitle

\begin{center}
\small
\texttt{galizur@gmail.com}
\end{center}

\begin{abstract}
We construct a relative category of exponential Nori motives over a smooth
algebraic base and prove an exponential analogue of the Nori--Ayoub functional
period theorem.  For a lisse relative exponential motive, the maximal trivial
subconnection of its exponential Gauss--Manin realization is induced by a
maximal constant exponential submotive.  Over $\overline{\Q}$, for every finitely generated
Tannakian subcategory of the generic relative Kontsevich--Zagier exponential
envelope and every $\lambda\in\overline{\Q}^{\times}$, we obtain an exact sequence
\[
  1\longrightarrow G_{\diff}
   \longrightarrow G_{\mot}
   \longrightarrow G_{\const}
   \longrightarrow 1.
\]
At $\lambda=1$, the generic relative formal exponential-period map is
injective after imposing the algebraic relations among constant exponential
periods.

The proof combines Sabbah's relative Fourier--nearby-cycle comparison with the
arithmetic theory of $1$-motives of Huber--W\"ustholz to control differential
characters.  Numerical relations arising only after specialization, and an
elementary presentation of all motivic period relations, are separate
questions.
\end{abstract}

\tableofcontents

\section{Introduction}

\subsection{The functional problem}

For the arithmetic results we work over
\[
  k=\overline{\Q},
\]
and let $S$ be a smooth connected quasi-projective $k$-variety.  The
exponentiated category and the fixed-part theorem are defined over any
algebraically closed subfield $k\subset\C$.  We use the $k$-linear scalar
extension of the rational Nori categories.

A relative exponential period is represented schematically by a family
\[
  s\longmapsto \int_{\gamma_s} e^{-f_s}\omega_s,
\]
where the de Rham class is computed with the twisted differential
$d-df_s\wedge$ and the Betti class is a rapid-decay cycle.  Exponential periods
admit equivalent naive, cohomological, and Nori-motivic descriptions
\cite{CHH,CommelinHuberII,FresanJossenExp}.  In families, the resulting functions
are solutions of exponential Gauss--Manin systems and admit rapid-decay integral
representations \cite{HienRoucairol,Hien}.

For ordinary periods, Nori and Ayoub prove a functional form of the
Kontsevich--Zagier period conjecture: the relative formal period algebra injects
into the corresponding analytic function field.  Equivalently, the relative
motivic comparison torsor is governed by algebraic monodromy
\cite{KontsevichZagier,AyoubPeriods,AyoubRelative,Mostaed,BakkerTsimerman}.
The aim of this paper is to establish the analogous statement for exponential
motives.

Functional and numerical relations must be distinguished.  The functional
problem concerns identities valid on a dense open of the parameter space.  The
numerical problem asks whether additional relations appear after specialization
to an algebraic point.  The latter requires separate arithmetic input and is not
addressed here.

\subsection{Main results}

Let $\MN(X)$ be the abelian category of perverse Nori motives on a
quasi-projective $k$-variety $X$ \cite{IvorraMorel,Terenzi,Tubach}.  For
\[
  \pi_X:X\times\Aone\longrightarrow X,
\]
define
\begin{equation}\label{eq:intro-EMN}
  \EMN(X)
  :=\MN(X\times\Aone)\big/\pi_X^*[1]\MN(X).
\end{equation}

\begin{theorem}[Exponentiated perverse-Nori heart]\label{thm:intro-heart}
The essential image of
\[
  \pi_X^*[1]:\MN(X)\longrightarrow\MN(X\times\Aone)
\]
is a Serre subcategory.  The quotient in \eqref{eq:intro-EMN} is the heart of
the exponentiated derived coefficient system associated with $D^b(\MN(-))$.
It carries the six operations and the additive-convolution tensor product.  The
Nori-to-Hodge realization descends to an exact faithful functor
\[
  R_H^{\exp}:\EMN(X)\longrightarrow\EMH(X)
\]
compatible with these operations and with every nonzero exponential
realization.
\end{theorem}

Fix $\lambda\in\C^\times$.  Let $S$ have dimension $d$, let
$a:S\to\Spec k$ be the structure morphism, and let $N\in\EMN(S)$ be
$\lambda$-lisse.  Set
\begin{equation}\label{eq:intro-fixed}
  N^{\mathrm{fix}}
  :=\im\!\left(
       a^*\,{}^p\!H^{-d}(a_*N)[d]\longrightarrow N
     \right).
\end{equation}

\begin{theorem}[Invariant exponential fixed part]\label{thm:intro-fixed}
The $\lambda$-realization of $N^{\mathrm{fix}}$ is the maximal trivial
differential submodule of $\rho_\lambda(N)$.  The object $N^{\mathrm{fix}}$ is
the maximal constant subobject of $N$.
\end{theorem}

Let $\eta=\Spec k(S)$ and let $\Tcal_{\mathrm{KZ}}(\eta)$ denote the generic
rigid Tannakian envelope generated by elementary relative
Kontsevich--Zagier exponential data.

\begin{theorem}[Relative exponential Galois exactness]
\label{thm:intro-exact}
Let $\lambda\in\overline{\Q}^{\times}$ and let $\Tcal$ be a finitely generated
rigid Tannakian subcategory of $\Tcal_{\mathrm{KZ}}(\eta)$.  Then the relative
differential Galois group is normal in the motivic Galois group, and
\[
  1\longrightarrow G_{\diff}
   \longrightarrow G_{\mot}
   \longrightarrow G_{\const}
   \longrightarrow1
\]
is exact.
\end{theorem}

\begin{theorem}[Relative exponential Nori--Ayoub theorem]
\label{thm:intro-functional}
For the canonical realization $\lambda=1$, after the algebraic relations
among constant exponential periods have been imposed, the generic relative
formal exponential-period map of the filtered Kontsevich--Zagier Tannakian
envelope is injective.  The same statement holds for every
$\lambda\in\overline{\Q}^{\times}$.
\end{theorem}

The fixed-part theorem follows from the lowest algebraic direct image.  The proof
of \cref{thm:intro-exact} requires control of differential characters.
Sabbah's relative Fourier--nearby-cycle comparison places their Betti local
systems in the ordinary geometric-origin category; Lam--Litt then gives finite
monodromy, and Huber--W\"ustholz's theory of $1$-motives produces the required
exponential-motivic lifts.  The resulting semi-invariant subobjects yield
normality by Andr\'e's criterion, and the comparison torsor is identified with
the Picard--Vessiot torsor.

\subsection{Organization}

\Cref{sec:background} reviews the coefficient systems and realizations.
\Cref{sec:Virk} records two consequences of Virk's realization theory.
The relative exponential Nori heart is constructed in \cref{sec:Nori-heart},
and the fixed-part theorem is proved in \cref{sec:fixed}.  Differential
characters are treated in \cref{sec:rank-one}.  The Galois exact sequence and
functional period theorem appear in \cref{sec:Tannaka,sec:periods}, and
\cref{sec:KZ} states the Kontsevich--Zagier consequence.  The examples in
\cref{sec:examples} include a complete functional presentation for finite
cyclotomic Gamma packets and isolate the classical Rohrlich--Lang problem in
the constant Gamma values.

\section{Background and conventions}\label{sec:background}

\subsection{Perverse Nori motives}

Let $\MNQ(X)$ be the rational category of perverse Nori motives constructed by
Ivorra--Morel.  Throughout the paper we use its scalar extension
\[
  \MN(X):=\MNQ(X)\otimes_{\Q}k,
\]
and similarly scalar-extend the Hodge and exponential realization categories; we
suppress the coefficient field from the notation.  All functors used below are
$\Q$-linear and extend $k$-linearly; on lisse Tannakian subcategories this is
the usual scalar extension \cite{DeligneMilne}.  We take $k$ algebraically
closed so that finite-order rank-one characters and their idempotents are
defined over $k$.  The Betti realization becomes a faithful exact functor
\[
  \operatorname{rat}_X:\MN(X)\longrightarrow
  \Perv(X^{\mathrm{an}},k).
\]
Its bounded derived categories carry the four operations and nearby and vanishing
cycles \cite{IvorraMorel}.  Terenzi constructs a canonical closed symmetric
monoidal structure and completes the six-functor formalism
\cite{Terenzi}.  Tubach proves that the corresponding derived categories are the
derived categories of their constructible hearts and constructs operation-compatible
Nori and Hodge realizations \cite{Tubach}.  We use the following consequences.

\begin{enumerate}[label=\textup{(N\arabic*)},leftmargin=3em]
\item $\operatorname{rat}_X$ is faithful and exact, and its derived extension is
conservative.
\item Smooth pullback with the conventional relative-dimension shift is
$t$-exact.
\item The six operations, tensor product, duality, and perverse cohomology commute
with Betti realization.
\item There is a faithful exact Hodge realization
\[
  R_{H,X}:\MN(X)\longrightarrow\MHM(X)
\]
whose derived extension commutes with the six operations.
\end{enumerate}

The categorical construction and fixed-part theorem descend from the rational
category by exact scalar extension.  The arithmetic exponential-kernel theorem
and the functional period theorem are stated over $\overline{\Q}$.

\subsection{Exponentiation}

Let $\Ccal$ be a six-functor coefficient system.  Gallauer--Pepin Lehalleur
associate with it an exponentiated coefficient system $\Ccal_{\exp}$ whose objects
over $X$ may be represented on $X\times\Ga$, modulo objects pulled back from $X$.
The tensor product is additive convolution in the potential coordinate
\cite{GallauerPepin}.  In Virk's Hodge/$D$-module model one sets
\[
  E(X)=D(X\times\Ga)\big/\pi_X^*D(X),
\]
and the quotient inherits the six operations and convolution
\cite[Sections~3--6]{Virk}.

For each $\lambda\in\C^\times$, Virk defines a realization
\[
  \real_\lambda:E(X)\longrightarrow D_h^b(X)
\]
by twisting with the exponential kernel $e^{-\lambda t}$ and integrating in the
$t$-direction.  He proves that it commutes with the four operations and duality,
is $t$-exact, and is faithful on the exponential heart
\cite[Theorems~6.13, 7.5, and 7.6]{Virk}.

\subsection{Rapid decay and exponential periods}

For a smooth variety $X$, a regular function $f$, and an algebraic differential
form $\omega$, the twisted de Rham complex uses $d-df\wedge$.  Its topological
dual is rapid-decay homology.  Hien proves the perfect period pairing for flat
algebraic connections, and Hien--Roucairol identify solutions of exponential
Gauss--Manin systems with rapid-decay integrals
\cite{Hien,HienRoucairol}.  Fresan--Jossen's exponential Nori category packages
the de Rham, rapid-decay, and perverse realizations
\cite{FresanJossenExp}.  Commelin--Habegger--Huber and
Commelin--Huber compare the naive, cohomological, and motivic notions of
exponential period \cite{CHH,CommelinHuberII}.

\begin{proposition}[Absolute rapid-decay fiber functor]
\label{prop:absolute-rd-input}
For the absolute exponential Nori envelope over $k$, the perverse realization
lands in Katz's convolution category $\Perv_0(\Aone)$.  The nearby-fiber functor
at infinity
\[
  \Psi_\infty:\Perv_0(\Aone)\longrightarrow\operatorname{Vec}_k
\]
is exact, faithful, and symmetric monoidal.  On an elementary exponential pair,
its value is canonically identified with rapid-decay Betti cohomology, compatibly
with the twisted de Rham comparison.
\end{proposition}

\begin{proof}
This combines Katz's convolution theorem with the perverse realization of
Fres\'an--Jossen and its comparison with rapid-decay cohomology; see
\cite{KatzESDE,FresanJossenExp} and
\cite[Corollary~2.17 and Theorem~3.6]{Snodgrass}.
\end{proof}

For a geometric point $s\in S(\C)$ and a lisse object $N\in\EMN(S)$, we use
the rapid-decay Betti fiber
\[
  \omega_{\Betti,s}(N)
  :=\Psi_\infty\!\left(
      \operatorname{rat}^{\exp}(i_s^*N)
    \right),
\]
where $i_s^*N$ is the absolute exponential Nori motive at $s$ and
$\Psi_\infty$ is the nearby-fiber-at-infinity fiber functor on the perverse
realization.  For elementary exponential pairs this is the usual rapid-decay
homology, and it is exact and monoidal on the lisse Tannakian subcategory
\cite{FresanJossenExp,Hien,Snodgrass}.  The de Rham/rapid-decay comparison
therefore supplies the comparison point used in \cref{sec:periods}.


\section{Consequences of Virk's realization theory}\label{sec:Virk}

The following consequences of Virk's realization theory separate the categorical and differential levels.

\begin{proposition}[Derived conservativity and reflection of the $t$-structure]
\label{prop:t-reflection}
Let $K\in E(X)$ and $\lambda\in\C^\times$.  Then
\[
  K\in E(X)^{\leq0}
  \Longleftrightarrow
  \real_\lambda(K)\in D_h^b(X)^{\leq0},
\]
and similarly for $\geq0$.  Thus $\real_\lambda$ is conservative on $E(X)$.
\end{proposition}

\begin{proof}
Virk proves $t$-exactness and faithfulness on the heart
\cite[Theorems~7.5--7.6]{Virk}.  Therefore
\[
  {}^pH^i\real_\lambda(K)
  \simeq \real_\lambda({}^pH^iK).
\]
If the left side vanishes, faithfulness on the heart gives ${}^pH^iK=0$.
Applying this degree by degree proves the assertions.
\end{proof}

\begin{proposition}[Exactness of the exponential Fourier transform]
\label{prop:FT-exact}
Let $V\to S$ be a vector bundle.  Virk's exponential Fourier transform
\[
  \mathsf{FT}_V:E(V)\longrightarrow E(V^\vee)
\]
is $t$-exact and induces an exact equivalence of exponential hearts.
\end{proposition}

\begin{proof}
Virk proves Fourier inversion and compatibility of $\mathsf{FT}_V$ with the
classical $D$-module Fourier transform under $\real_1$
\cite[Theorems~5.3 and 6.12]{Virk}.  With Virk's shifts, the classical Fourier
transform is $t$-exact.  Hence the realization of $\mathsf{FT}_V(K)$ lies in the
heart whenever $K$ does.  By \cref{prop:t-reflection},
$\mathsf{FT}_V(K)$ lies in the exponential heart.  Fourier inversion proves that
the induced exact functor on hearts is an equivalence.
\end{proof}

\section{The relative exponentiated Nori heart}\label{sec:Nori-heart}

Fix a quasi-projective $k$-variety $X$ and write
\[
  \pi=\pi_X:X\times\Aone\longrightarrow X.
\]
Since $\pi$ is smooth of relative dimension one, $\pi^*[1]$ is $t$-exact on
perverse Nori motives.

\begin{lemma}[Reflection of constant objects]\label{lem:constant-reflection}
Let $A\in\MN(X\times\Aone)$.  If
\[
  \operatorname{rat}(A)\simeq\pi^*L[1]
\]
for some $L\in\Perv(X^{\mathrm{an}},k)$, then there exists
$B\in\MN(X)$ such that $A\simeq\pi^*B[1]$.
\end{lemma}

\begin{proof}
Set
\[
  B:={}^{p}H^{-1}(\pi_*A).
\]
Compatibility with Betti realization and $\Aone$-homotopy invariance give
\[
  \operatorname{rat}(B)\simeq L.
\]
The adjunction counit yields a morphism
\[
  \pi^*B[1]\longrightarrow A.
\]
Its Betti realization is the canonical isomorphism
$\pi^*L[1]\to\pi^*L[1]$.  Since the derived Betti realization is conservative,
the counit is an isomorphism.
\end{proof}

\begin{proposition}[Serre property]\label{prop:Serre}
The essential image of
\[
  \pi^*[1]:\MN(X)\longrightarrow\MN(X\times\Aone)
\]
is closed under subobjects, quotients, and extensions.
\end{proposition}

\begin{proof}
For perverse sheaves, the essential image of smooth pullback along a surjective
smooth morphism with geometrically connected fibers is Serre
\cite[Corollaire~4.2.6.2]{BBD}; Virk applies this to $\pi$ in
\cite[Proposition~7.2]{Virk}.  Let $A$ be a subobject, quotient, or extension of
objects in the Nori essential image.  Exactness of Betti realization places
$\operatorname{rat}(A)$ in the corresponding Serre subcategory of perverse
sheaves.  The conclusion follows from \cref{lem:constant-reflection}.
\end{proof}

\begin{lemma}[Derived constant subcategory]\label{lem:derived-constant}
Let
\[
  \Ccal_X:=\pi^*[1]\MN(X)\subset\MN(X\times\Aone).
\]
The exact functor \(\pi^*[1]\) is fully faithful, and its derived functor
identifies \(D^b(\MN(X))\) with the full triangulated subcategory of
\(D^b(\MN(X\times\Aone))\) consisting of complexes whose perverse
cohomology objects lie in \(\Ccal_X\).  This subcategory is stable under
perverse truncations.
\end{lemma}

\begin{proof}
Full faithfulness follows from adjunction and \(\Aone\)-homotopy invariance.
By \cref{prop:Serre}, \(\Ccal_X\) is a Serre subcategory.  Tubach's
derived-category theorem identifies the constructible derived category with the
bounded derived category of its perverse-Nori heart.  Standard d'evissage for a
Serre subcategory then identifies \(D^b(\Ccal_X)\) with the full subcategory
of complexes whose perverse cohomology lies in \(\Ccal_X\); see
\cite[Section~1.4]{BBD}.  The subcategory is therefore stable under perverse truncations.  Finally,
\(\pi^*[1]:\MN(X)\to\Ccal_X\) is an equivalence, hence so are the
bounded derived categories.
\end{proof}

\begin{definition}[Relative exponential Nori heart]\label{def:EMN}
Define
\[
  \EMN(X)
  :=\MN(X\times\Aone)\big/\pi^*[1]\MN(X).
\]
\end{definition}

\begin{theorem}[Exponentiated perverse-Nori heart]\label{thm:heart-full}
The quotient in \cref{def:EMN} is the heart of
\[
  D^b(\MN(X\times\Aone))
  \big/\pi^*D^b(\MN(X)).
\]
As $X$ varies, these derived categories form the exponentiation of the
perverse-Nori coefficient system.  They carry the six operations and the
additive-convolution tensor product.  The Hodge realization descends to an exact
faithful operation-compatible functor
\[
  R_H^{\exp}:\EMN(X)\longrightarrow\EMH(X).
\]
\end{theorem}

\begin{proof}
By \cref{lem:derived-constant}, the thick subcategory being quotiented is stable
under perverse truncations.  The localization theorem for $t$-structures therefore
endows the Verdier quotient with a bounded $t$-structure whose heart is the Serre
quotient in \cref{def:EMN}; compare \cite[Section~1.4]{BBD} and Virk's proof in
\cite[Section~7]{Virk}.  Terenzi and Tubach provide a closed symmetric monoidal
six-functor coefficient system on $D^b(\MN(-))$ \cite{Terenzi,Tubach}.
Gallauer--Pepin Lehalleur's exponentiation theorem therefore supplies the six
operations and convolution on the quotient \cite{GallauerPepin}.

The Hodge realization commutes with pullback and carries
$\pi^*[1]\MN(X)$ into $\pi^*[1]\MHM(X)$, so it descends to the quotient.  It is
exact because the source and target quotient functors are exact.  To prove
faithfulness, let $f:A\to B$ have zero image in the Hodge quotient.  Then
$R_H(\im f)$ belongs to the Hodge constant Serre subcategory.  Its underlying
perverse sheaf is pulled back from $X$, so \cref{lem:constant-reflection} gives
$\im f\in\pi^*[1]\MN(X)$.  Hence $f$ was already zero in the Nori quotient.
Compatibility with the operations is inherited from Tubach's realization and
the functoriality of exponentiation.
\end{proof}

For $\lambda\in\C^\times$ put
\[
  \rho_\lambda:=\real_\lambda\circ R_H^{\exp}.
\]

\begin{definition}[Lisse and admissible Tannakian subcategories]
\label{def:lisse-objects}
An object $N\in\EMN(X)$ is \emph{$\lambda$-lisse} if
$\rho_\lambda(N)$ is a vector bundle with integrable connection.  The full
$\lambda$-lisse subcategory is a Serre tensor subcategory of $\EMN(X)$,
because $\rho_\lambda$ is exact, faithful, and tensor-compatible and algebraic
connections form a Serre tensor subcategory of holonomic $\mathscr D$-modules.

A \emph{$\lambda$-admissible Tannakian subcategory} is a strictly full rigid
tensor subcategory of the $\lambda$-lisse heart which is closed under
subquotients.  All Tannakian arguments below take place in such a subcategory.
When the scale is fixed, we suppress $\lambda$.
\end{definition}

\begin{remark}[Absolute exponential motives]\label{rem:absolute-FJ}
For $X=\Spec k$, recollement identifies the quotient with the Nori envelope
generated by perverse objects on $\Aone$ with vanishing global cohomology, the
model used by Fres\'an--Jossen for absolute exponential motives
\cite{FresanJossenExp}.
\end{remark}

\section{The invariant fixed-part theorem}\label{sec:fixed}

Let $S$ be smooth, connected, and of dimension $d$, let
$a:S\to\Spec k$ be the structure morphism, and fix a scale
$\lambda\in\C^\times$.  Throughout this section $N\in\EMN(S)$ is
$\lambda$-lisse.

\begin{definition}[Constant subcategory]
The constant exponential subcategory of $\EMN(S)$ is the essential image of
\[
  a^*[d]:\EMN(\Spec k)\longrightarrow\EMN(S).
\]
\end{definition}

\begin{lemma}[Lowest algebraic direct image]\label{lem:lowest-direct-image}
Let $M$ be a vector bundle with integrable algebraic connection on the smooth
$d$-fold $S$.  In the perverse normalization of algebraic $\mathscr D$-modules,
\[
  {}^p\!H^j(a_+M)=0\quad(j<-d),
  \qquad
  {}^p\!H^{-d}(a_+M)=H^0_{\dR}(S,M)=H^0_\nabla(S,M).
\]
\end{lemma}

\begin{proof}
The algebraic de Rham complex of $M$ is concentrated in ordinary degrees
$0,\ldots,d$.  The perverse normalization shifts a connection by $d$, giving the stated
cohomological amplitude and identifying the lowest group with algebraic
horizontal sections; see \cite[Chapters~1 and~7]{HTT}.
\end{proof}

Put
\[
  C_N:={}^{p}H^{-d}(a_*N)\in\EMN(\Spec k).
\]
By operation compatibility and \cref{lem:lowest-direct-image}, the perverse
cohomology of $a_*N$ vanishes below degree $-d$ after applying
$\rho_\lambda$.  Faithfulness reflects this vanishing to $\EMN$.  Hence the
canonical truncation map $C_N[d]\to a_*N$ and adjunction give
\[
  \epsilon_N:a^*C_N[d]\longrightarrow N.
\]
Define
\[
  N^{\mathrm{fix}}:=\im(\epsilon_N).
\]

\begin{lemma}[Horizontal evaluation]\label{lem:horizontal-evaluation}
Let $M=\rho_\lambda(N)$.  The realization of $\epsilon_N$ is the evaluation
map
\[
  \Ocal_S\otimes_k H^0_\nabla(S,M)\longrightarrow M.
\]
Its image is the maximal trivial differential submodule of $M$.
\end{lemma}

\begin{proof}
By \cref{lem:lowest-direct-image}, the realization of $C_N$ is
$H^0_\nabla(S,M)$, and the adjunction morphism realizes as evaluation.  Its
image is trivial.  Conversely, every trivial differential submodule has a
basis of global algebraic horizontal sections and is therefore contained in
the evaluation image.
\end{proof}

\begin{lemma}[Constant reflection]\label{lem:constant-reflection-exp}
Let $Q\in\EMN(S)$ be $\lambda$-lisse.  If $\rho_\lambda(Q)$ is a trivial
connection, then the adjunction evaluation
\[
  a^*\,{}^pH^{-d}(a_*Q)[d]\longrightarrow Q
\]
is an isomorphism.  Thus $Q$ is constant.
\end{lemma}

\begin{proof}
After realization the displayed map is the evaluation isomorphism
\[
  \Ocal_S\otimes_k H^0_\nabla(S,\rho_\lambda(Q))
  \xrightarrow{\sim}\rho_\lambda(Q).
\]
Exact faithfulness of $\rho_\lambda$ reflects this isomorphism.
\end{proof}

\begin{proposition}[Subquotient closure of constants]
\label{prop:constant-subquotients}
Inside the $\lambda$-lisse heart, the constant subcategory is closed under
subobjects and quotients, as well as tensor products and finite direct sums.
\end{proposition}

\begin{proof}
A subobject or quotient of a trivial connection is trivial.  Exactness of
$\rho_\lambda$ and \cref{lem:constant-reflection-exp} therefore show that every
subobject or quotient of a constant $\lambda$-lisse object is constant.  Smooth pullback preserves tensor products and finite direct sums.
\end{proof}

\begin{theorem}[Invariant exponential fixed part]\label{thm:fixed-part}
For the fixed scale $\lambda$,
\[
  \rho_\lambda(N^{\mathrm{fix}})
  =\Fix\!\left(\rho_\lambda(N)\right).
\]
Moreover, $N^{\mathrm{fix}}$ is constant and is the maximal constant
exponential Nori subobject of $N$.
\end{theorem}

\begin{proof}
Exactness carries the image defining $N^{\mathrm{fix}}$ to the evaluation
image, so the first assertion is \cref{lem:horizontal-evaluation}.  This
realization is trivial; hence \cref{lem:constant-reflection-exp} shows that
$N^{\mathrm{fix}}$ is constant.

Let $a^*B[d]\hookrightarrow N$ be any constant subobject.  By adjunction it
corresponds to a morphism $B[d]\to a_*N$.  Since the perverse cohomology of
$a_*N$ vanishes below degree $-d$, the morphism factors through
$C_N[d]\to a_*N$.  Its adjoint factors through $\epsilon_N$, so the subobject
is contained in $N^{\mathrm{fix}}$.
\end{proof}

\section{Differential characters and Fourier geometry}\label{sec:rank-one}

\subsection{Rank-one lifting}

Fix a nonzero scale $\lambda$ and recall
$\rho_\lambda=\real_\lambda\circ R_H^{\exp}$ from
\cref{def:lisse-objects}.  Let $\Tcal\subset\EMN(S)$ be a rigid Tannakian
subcategory of $\lambda$-lisse objects, and
let $\Tcal_\nabla$ be the Tannakian subcategory of algebraic connections generated
by $\rho_\lambda(\Tcal)$ and all its differential subquotients.

\begin{definition}[Rank-one saturation]\label{def:RS}
The pair $(\Tcal,\rho_\lambda)$ is rank-one saturated if for every invertible
$L\in\Tcal_\nabla$ there is an invertible object $\mathcal L$ in the ambient
relative exponential Nori category such that
\[
  \rho_\lambda(\mathcal L)\simeq L.
\]
\end{definition}

The lift may lie in the ambient category; the associated semi-invariant is
recovered inside an object of $\Tcal$ in \cref{sec:Tannaka}.

\subsection{A split Fourier-geometric criterion}

We use the following sufficient condition.

\begin{definition}[Split Fourier-geometric category]\label{def:split-FG}
We call $(\Tcal,\rho_\lambda)$ \emph{split Fourier-geometric} if, for every
$B\in\Tcal_\nabla$, there exist a dense open $S^\circ\subset S$, a finite Galois
cover $q:S'\to S^\circ$, and a good compactification such that:
\begin{enumerate}[label=\textup{(FG\arabic*)},leftmargin=3.5em]
\item the Stokes-graded connection admits a global decomposition
\[
  \gr^{\St}(q^*B)
   \simeq\bigoplus_{\alpha\in A_B}
      \Ecal^{g_\alpha}\otimes R_\alpha,
\]
with $g_\alpha\in\Ocal(S')$ pairwise distinct modulo constants and with
$R_\alpha$ regular singular;
\item every $R_\alpha$ is of geometric origin in the ordinary Nori category;
\item the Galois descent datum on every rank-one summand
$\Ecal^{g_\alpha}\otimes A$ with $A$ of finite order is effective in the
relative exponential Nori category.
\end{enumerate}
\end{definition}

\begin{lemma}[Rank-one selection of a Stokes grade]\label{lem:one-grade}
Let $L$ be a rank-one differential subquotient of a connection satisfying
\textup{(FG1)}.  Then, after the same finite cover,
\[
  q^*L\simeq\Ecal^{g_\alpha}\otimes A
\]
for one $\alpha$, where $A$ is a rank-one regular-singular subquotient of
$R_\alpha$.
\end{lemma}

\begin{proof}
Good formal decomposition and the Stokes-graded functor are exact in the abelian
category of Stokes-filtered local systems.  A rank-one connection has one formal
exponential factor and no off-diagonal Stokes operator.  Therefore its graded
object is a rank-one subquotient of one summand.  Removing the selected
exponential factor leaves a rank-one regular-singular subquotient of
$R_\alpha$.
\end{proof}

\begin{lemma}[Finite monodromy of the regular factor]\label{lem:regular-finite}
Under \textup{(FG2)}, the rank-one regular factor $A$ in
\cref{lem:one-grade} has finite monodromy.
\end{lemma}

\begin{proof}
Subquotients of ordinary geometric local systems are again of geometric origin.
Lam--Litt prove that a rank-one complex local system on a smooth variety is of
geometric origin if and only if it is torsion
\cite[Lemma~4.2.1]{LamLitt}.  Hence $A$ has finite monodromy.  Over the algebraically closed coefficient
field, the corresponding character idempotent in the finite \'etale
pushforward defines an invertible Artin object.
\end{proof}

\begin{theorem}[Generic split Fourier-geometric saturation]
\label{thm:split-saturation}
Let $\Tcal$ be a finitely generated split Fourier-geometric rigid tensor
subcategory.  After restriction to a suitable common dense open
$S^\circ\subset S$, the pair $(\Tcal|_{S^\circ},\rho_\lambda)$ is rank-one
saturated.
\end{theorem}

\begin{proof}
Let $L\in\Tcal_\nabla$ be invertible.  By
\cref{lem:one-grade,lem:regular-finite}, after a finite Galois cover
\[
  q^*L\simeq\Ecal^{g}\otimes A
\]
with $g$ algebraic and $A$ finite.  At the fixed scale $\lambda$, the graph
object with potential $-g/\lambda$ realizes $\Ecal^g$ under the convention
$\Ecal^g=(\Ocal,d-dg)$ and is invertible in the relative exponential Nori
category.  If $m$ is the order of $A$, choose
the finite \'etale cover which trivializes $A$.  Over $k$, the character idempotent attached to $A$ cuts out an
invertible rank-one Artin object.  The convolution tensor product of this Artin object with the graph object
therefore realizes $q^*L$.  Hypothesis \textup{(FG3)} then produces an invertible object over $S^\circ$ realizing $L|_{S^\circ}$.  Since
$\Tcal$ is finitely generated, the finitely many dense opens required for its
generators and tensor constructions may be replaced by one common dense open.
Hence the restricted category is rank-one saturated.
\end{proof}

\subsection{Nonzero Fourier fibers and nearby cycles}

Let $K$ be a regular geometric perverse-Nori object on
$S\times\Aone_t$, and let
\[
  \widehat K=\operatorname{FL}_{t/S}(K)
\]
be its partial Fourier--Laplace transform on
$S\times\Aone_\lambda$.  Fix $\lambda_0\in\C^\times$ and shrink $S$ so that
\[
  M_{\lambda_0}:=i_{\lambda_0}^{+}\widehat K
\]
is a connection on $S$.

\begin{lemma}[Finite graph reduction]\label{lem:finite-graph}
After shrinking $S$, there is a finite \`etale cover
\[
  r:S'\longrightarrow S
\]
and pairwise disjoint smooth algebraic graphs
\[
  H_i=\{t=u_i(s)\}\subset S'\times\Aone_t,
  \qquad u_i\in\Ocal(S'),
\]
such that the characteristic variety of the regular holonomic
$\mathscr D$-module underlying $r^*K$ is contained in the union of the zero
section and the conormal bundles $T^*_{H_i}(S'\times\Aone_t)$.
\end{lemma}

\begin{proof}
Choose an algebraic Whitney stratification of $S\times\Aone_t$ adapted to the
Betti realization of $K$.  After deleting the images in $S$ of all strata which
do not dominate the base, every remaining non-open stratum has dimension
$\dim S$: the ambient variety has dimension $\dim S+1$, while a proper stratum
dominating $S$ has dimension at least $\dim S$ and codimension at least one.
It is therefore generically finite over $S$.  Normalize the finitely many such
multisections, shrink once more so that they are finite \`etale, and pass to a
common finite Galois refinement.  They then split into smooth sections.  Delete
from $S$ the Galois-stable union of the images of all pairwise intersections;
after pulling back this open, the graphs are pairwise disjoint.  The characteristic
variety of a regular holonomic module constructible for the chosen Whitney
stratification is contained in the union of the conormal varieties of the
strata.  The open stratum contributes the zero section, and the remaining
strata contribute the conormals of the displayed graphs.
\end{proof}

\begin{lemma}[Nonzero-scale lissity]\label{lem:nonzero-lisse}
After shrinking $S$ further,
\[
  \widehat K\big|_{S\times\Gm}
\]
is a connection.
\end{lemma}

\begin{proof}
Apply \cref{lem:finite-graph} and work analytically on coordinate polydiscs of
$S'$.  Put $z=\lambda^{-1}$ on the punctured Fourier line.  In the notation of
Sabbah's parameterized partial Laplace transform, the characteristic-variety
hypothesis of Appendix~B, Proposition~B.1 is now satisfied by the hypersurfaces
$H_i=\{t-u_i=0\}$.  Sabbah states on p.~40, immediately before Proposition~B.1, that the
localized transform is a locally free $\Ocal_{S'}[z,z^{-1}]$-module of finite
rank with an integrable connection.  Proposition~B.1 then identifies its formal
logarithmic residual systems with the vanishing-cycle connections along the
$H_i$.  The microlocalization and decomposition inputs are the results of
Douai--Sabbah cited there
\cite[Appendix~B, p.~40 and Proposition~B.1]{SabbahParameters}
\cite[Propositions~1.18 and~1.20]{DouaiSabbah}.  Under
$z=\lambda^{-1}$ this localized transform is
$r^*\widehat K|_{S'\times\Gm}$.  Hence the latter is lisse.  Lissity is local for
the finite \`etale topology, so it descends from $S'$ to $S$.
\end{proof}

By \cref{lem:nonzero-lisse}, the Betti realization of $\widehat K$ on
$S\times\Gm$ is, up to the standard perverse shift, a local system
$\mathscr G$.  It corresponds to a representation of
\[
  \pi_1(S\times\Gm)
  \simeq\pi_1(S)\times\mathbb Z.
\]
We use perverse-normalized cycles
\[
  {}^p\!\psi_f:=R\psi_f[-1],
  \qquad
  {}^p\!\phi_f:=R\phi_f[-1].
\]
Let $\mathscr G_\psi$ denote the unshifted local system on $S$ underlying
${}^p\!\psi_\lambda\operatorname{rat}(\widehat K)$, equipped with its commuting
nearby-monodromy operator $T_\lambda$.

\begin{proposition}[Fixed fiber versus nearby cycles]\label{prop:fiber-nearby}
After forgetting $T_\lambda$, the local systems
\[
  \mathscr G_{\lambda_0}
  \quad\text{and}\quad
  \operatorname{For}_{T_\lambda}(\mathscr G_\psi)
\]
are isomorphic, though not canonically so.  Compatible basepoints and a path in
$\Gm$ from $\lambda_0$ to a punctured neighbourhood of the origin determine
such an isomorphism; two choices differ by a power of $T_\lambda$.  Consequently, the fixed nonzero fiber and the nearby-cycle object have the
same Jordan--H\"older constituents as $\pi_1(S)$-representations.
\end{proposition}

\begin{proof}
By \cref{lem:nonzero-lisse}, $\mathscr G$ is a local system on
$S\times\Gm$.  Such a local system is a representation of
$\pi_1(S)\times\mathbb Z$.  Restriction to $S\times\{\lambda_0\}$ forgets the
central $\mathbb Z$-operator.  After choosing compatible basepoints and a path
to a punctured neighbourhood of the origin, perverse nearby cycles identify
the same vector space and the same $\pi_1(S)$-action, while retaining the
commuting $\mathbb Z$-operator as $T_\lambda$.  Changing the path winds around
the origin and composes the identification with a power of $T_\lambda$.  Thus
the underlying $S$-local systems are isomorphic, though not canonically so.
\end{proof}

\subsection{Sabbah's conventions and the relative comparison}

Compactify the source variable by writing
\[
  X=S\times\mathbb P^1_t,
  \qquad t'=t^{-1}
\]
near infinity, and let
\[
  q:X\times\Aone_\lambda\longrightarrow S\times\Aone_\lambda
\]
be the projective projection.  Let $\Mcal$ be the regular holonomic
$\mathscr D$-module underlying the Hodge realization of $K$ on
$S\times\Aone_t$, and let
\[
  \Tcal:=\Mcal_{\min}
\]
be its regular minimal (middle) extension across $t=\infty$.  Sabbah's Fourier kernel
uses the localization $\Tcal[*\infty]=\Mcal[*\infty]$.  Hence replacing an
arbitrary regular extension by the minimal extension does not change the
localized kernel, its partial Fourier--Laplace transform, or any nonzero
Fourier fiber under study.  He first constructs the Fourier kernel object
$\mathcal F\Tcal$ on $X\times\Aone_\lambda$ and only then applies $q_+$.
Because the Fourier kernel depends only on $\Tcal[*\infty]$, the normalized
classical partial Fourier--Laplace transform is the usual exact transform of the
localized regular holonomic perverse module.  Its projective direct image is
therefore concentrated in perverse degree zero:
\[
  {}^p\!\mathcal H^j q_+(\mathcal F\Tcal)=0\quad(j\ne0),
  \qquad
  \widehat{\Tcal}:={}^p\!\mathcal H^0q_+(\mathcal F\Tcal),
\]
and $\widehat{\Tcal}$ underlies $\widehat K$; see
\cite[Appendix~A]{Virk} and \cite[Chapter~3]{HTT}.  The unipotent case is treated before projective direct image.

We keep Sabbah's nearby-cycle exponent labels and write his Fourier coordinate
$\tau$ as $\lambda$.  With the perverse normalization above, his comparison has
the following form on underlying regular $\mathscr D$-modules and local systems:
\begin{equation}\label{eq:sabbah-comparison}
\begin{array}{c|c}
\text{label} & \text{comparison before }q_+ \\
\hline
-1\leq\operatorname{Re}\alpha<0,\ \alpha\neq-1
&
\Psi_{\lambda,\alpha}(\mathcal F\mathcal T)
 \simeq
i_{\infty,+}\Psi_{t',\alpha}(\mathcal T)
\\[1mm]
0
&
\phi_{\lambda,0}(\mathcal F\mathcal T)
 \simeq
i_{\infty,+}\psi_{t',-1}(\mathcal T)
\\[1mm]
-1
&
0\to i_{\infty,+}\ker N_{t'}
 \to\ker N_\lambda\to\mathcal T\to0.
\end{array}
\end{equation}
The first two lines are Proposition~5.8; the first two comparisons after
projective direct image are Corollary~5.20
\cite[Proposition~5.8 and Corollary~5.20]{SabbahMonodromy}.  The last exact
sequence is also part of Proposition~5.8.  Here $\phi_{\lambda,0}$ is the
vanishing-cycle object for eigenvalue $1$, whereas the unipotent nearby-cycle
object is $\psi_{\lambda,-1}$.

Two twistor normalizations in Sabbah's formulas do not alter the underlying
$\mathscr D$-module or local system.  By Sabbah's Lemma~5.6, the factor
$\Ocal_{\Omega_0}(-D_\alpha)$ is absorbed at the level of $R$-triples in the
comparison above.  The half twist $(-1/2)$ rescales the twistor pairing and
weight normalization.  Neither is a Nori half-Tate twist, so both are omitted
when discussing simple local-system constituents
\cite[Lemma~5.6]{SabbahMonodromy}.

\begin{theorem}[Nearby-cycle geometricity]\label{thm:nearby-geometric}
Every simple lisse constituent of
\[
  {}^p\!\psi_{\lambda=0}(\widehat K)
\]
is of ordinary geometric origin.
\end{theorem}

\begin{proof}
We first treat a simple pure constituent of the Hodge realization of $K$.
Its regular minimal extension $\Tcal=\Mcal_{\min}$ underlies a
polarizable regular twistor $\mathscr D$-module by the Hodge-to-twistor
realization \cite[Section~5]{SabbahTwistor}.  Sabbah's relative Fourier
comparison then applies with parameter space $Y=S$ \cite{SabbahMonodromy}.

For the non-unipotent labels, the first line of
\eqref{eq:sabbah-comparison}, followed by the corresponding comparison after
$q_+$, identifies every simple constituent of the Fourier nearby-cycle object
with a constituent of an ordinary nearby-cycle object at $t=\infty$.  The
second line treats the eigenvalue-one vanishing-cycle object.

For the unipotent nearby-cycle object
\[
  V={}^p\!\psi_{\lambda,-1}(\mathcal F\Tcal),
  \qquad N=N_\lambda,
\]
use the finite filtration
\[
  0\subset\ker N\subset\ker N^2\subset\cdots\subset V.
\]
For each $j\geq1$, the map induced by $N^{j-1}$ embeds
$\ker N^j/\ker N^{j-1}$ into $\ker N$.  Hence every simple constituent of $V$
occurs in $\ker N$.  Sabbah's pre-$q_+$ exact sequence expresses $\ker N$ as
an extension of $\Tcal$ by
$i_{\infty,+}\ker N_{t'}$.  Therefore every simple constituent of $V$ occurs
in
\begin{equation}\label{eq:unipotent-geometric-sources}
  \Tcal
  \quad\text{or}\quad
  i_{\infty,+}{}^p\!\psi_{t',-1}(\Tcal).
\end{equation}

Now apply the derived proper direct image $Rq_*$.  Proper compatibility of
nearby cycles identifies $Rq_*V$ with the unipotent nearby-cycle complex of
$q_+(\mathcal F\Tcal)$.  Applying $Rq_*$ to a finite constituent filtration of
$V$ gives distinguished triangles whose long exact perverse-cohomology
sequences show that every simple constituent of
${}^p\!\psi_{\lambda,-1}(\widehat K)$ occurs among the perverse-cohomology
constituents of the proper direct images of the two ordinary geometric objects
in \eqref{eq:unipotent-geometric-sources}.  Thus Sabbah's unipotent sequence is used only before projective direct image.

Middle extension, nearby cycles, proper direct image, and perverse cohomology
are operations in the ordinary perverse-Nori coefficient system
\cite{IvorraMorel,Tubach}.  Hence every constituent produced above is of
ordinary Nori-geometric origin.

For mixed $K$, work on its mixed Hodge realization
\[
  H:=R_H(K)\in\MHM(S\times\Aone_t).
\]
Saito's weight filtration is strict, and normalized partial Fourier transform
and perverse nearby cycles are exact.  They therefore induce a finite filtration
on ${}^p\!\psi_\lambda(\operatorname{FL}_{t/S}H)$ whose graded subquotients
are subquotients of
\[
  {}^p\!\psi_\lambda
  \operatorname{FL}_{t/S}(\operatorname{gr}^W_mH).
\]
Each pure graded object is semisimple.  Its simple strict-support constituents
belong, on the underlying local-system side, to the abelian subquotient-closed
category of local systems of geometric origin identified by
Jacobsen--Terenzi \cite[Corollary~5.15]{JacobsenComparison}.  The pure argument
above applies constituent by constituent.  Hence every simple constituent of
the mixed nearby-cycle object is of geometric origin \cite{SaitoMHM}.  Sabbah's comparison is used only to identify complex isomorphism classes;
geometric origin is supplied by the ordinary geometric objects on the right.
\end{proof}

\begin{remark}[The two-graph family]\label{rem:two-graph}
For the family $y^2=t(t-s)$, the local system extends at $t=\infty$, so
${}^p\!\psi_{t',-1}\Tcal$ has rank one, $N_{t'}=0$, and trivial
$s$-monodromy.  Moreover $H^0(q_+\Tcal)=H^2(q_+\Tcal)=0$ because there are no
invariant or coinvariant vectors, while
$\chi(\mathbb P^1\setminus\{0,s\},\mathcal L)=0$; hence
$q_+\Tcal=0$.  Sabbah's upstream exact sequence gives
$\dim\ker N_\lambda=1$, whereas the unipotent nearby-cycle object has rank
two.  Thus $N_\lambda$ has a single Jordan block and the two rank-one
Jordan--H\"older constituents are trivial, as in the confluent system
$M(1/2,1,\lambda s)$.
\end{remark}

\begin{corollary}[Finite monodromy of rank-one Fourier factors]
\label{cor:rankone-finite}
Every rank-one Jordan--H\"older constituent of a fixed nonzero Fourier fiber has
finite Betti monodromy.
\end{corollary}

\begin{proof}
By \cref{prop:fiber-nearby}, after forgetting the nearby-monodromy
operator the fixed fiber and the nearby-cycle object have the same underlying
$\pi_1(S)$-representation.  By \cref{thm:nearby-geometric}, the resulting
rank-one constituent lies in the Nori category of local systems of geometric
origin.  Jacobsen--Terenzi identify this category with Ayoub's category of local
systems of geometric origin \cite[Corollary~5.15]{JacobsenComparison}.
Lam--Litt prove that a
rank-one complex local system on a smooth variety is of geometric origin if and
only if it is torsion \cite[Lemma~4.2.1]{LamLitt}.
\end{proof}

\subsection{The arithmetic exponential kernel}

From this point $k=\overline{\Q}$, and we use the convention
$\Ecal^f=(\Ocal_U,d-df)$.

\begin{proposition}[Picard $1$-motive convention]
\label{prop:picard-convention}
Let $U$ be smooth and let $\overline U$ be a smooth proper compactification
with normal-crossings boundary $D$.  With the covariance conventions used in
this paper,
\[
  T_H\!\left(\operatorname{Pic}^{+}(U)\right)
  \simeq H^1(U,\mathbb Z(1)),
  \qquad
  \operatorname{Pic}^{+}(U)
  =\left[\operatorname{Div}^{0}_{D}(\overline U)
     \longrightarrow\operatorname{Pic}^{0}(\overline U)\right].
\]
A morphism $\Q(0)\to\operatorname{Pic}^{+}(U)_{\Q}$ is represented by a
rational boundary divisor in the kernel of the map to
$\operatorname{Pic}^{0}(\overline U)\otimes\Q$.  After multiplication by an
integer making the divisor principal, say $\operatorname{div}(u)$, its de Rham
realization is the logarithmic class $[d\log u]$.
\end{proposition}

\begin{proof}
This is the cohomological Picard $1$-motive and its Hodge and de Rham
realizations as constructed by Barbieri--Viale and Srinivas; see
\cite[Sections~2--4]{BarbieriVialeSrinivas} and
\cite[Section~10]{DeligneHodgeIII}.
\end{proof}

\begin{theorem}[Arithmetic exponential-kernel theorem]
\label{thm:arithmetic-kernel}
Let $U/\overline{\Q}$ be smooth and connected, and let $L$ be a rank-one
algebraic integrable connection over $\overline{\Q}$.  If its Betti local system
has finite monodromy, then there exist a dense affine open immersion
$j:U^\circ\hookrightarrow U$, a function $f\in\Ocal(U^\circ)$, and a
finite-order rank-one algebraic connection $A$ on $U^\circ$ such that
\[
  j^*L\simeq\Ecal^f\otimes A.
\]
\end{theorem}

\begin{proof}
Let $A_0$ be the finite Artin connection with the same Betti character as $L$;
it exists over $\overline{\Q}$ by invariance of finite \`etale covers under
extension of algebraically closed fields.  Set
$L_0=L\otimes A_0^{-1}$, so that $L_0$ has trivial Betti monodromy.  Choose a
dense affine open immersion $j:U^\circ\hookrightarrow U$ on which the
underlying line bundle of $L_0$ is trivial, and replace $U,L_0$ by
$U^\circ,j^*L_0$ for the remainder of the proof.  Write
\[
  \nabla=d+\omega,
  \qquad
  \omega\in\Gamma(U,\Omega^1_{U/\overline{\Q}}),
  \qquad d\omega=0.
\]
Trivial monodromy gives
\[
  \int_\gamma\omega\in2\pi i\mathbb Z
  \qquad(\gamma\in H_1(U(\C),\mathbb Z)).
\]
Let $\eta\in H^1_{\Betti}(U,\Q)$ be the integral class defined by
\[
  \langle\eta,\gamma\rangle
  =(2\pi i)^{-1}\int_\gamma\omega.
\]
Then comparison sends $[\omega]_{\dR}$ to $2\pi i\eta$.  If
$\eta=0$, comparison gives $[\omega]_{\dR}=0$; since $U$ is affine there is
$F\in\Ocal(U)$ with $\omega=dF$, and $L_0\simeq\Ecal^{-F}$.  We henceforth assume $\eta\neq0$.

The de Rham--Betti object $H^1_{\dRB}(U)$ then contains the one-dimensional
subobject determined by the specific pair
$([\omega]_{\dR},\eta_{\Betti})$:
\begin{equation}\label{eq:tate-line}
  \Q(-1)\hookrightarrow H^1_{\dRB}(U),
\end{equation}
where $\Q(-1)$ is the Tate object whose period is $2\pi i$.  Twisting by
$\Q(1)$ gives
\begin{equation}\label{eq:tate-twist}
  \Q(0)\hookrightarrow H^1_{\dRB}(U)(1).
\end{equation}

Ayoub--Barbieri-Viale identify the degree-one Nori category with Deligne
$1$-motives with torsion \cite{AyoubBarbieriViale}.  By \cref{prop:picard-convention}, the cohomological Picard $1$-motive of
$U$ is
\begin{equation}\label{eq:pic-plus}
  \operatorname{Pic}^{+}(U)
  =\left[
      \operatorname{Div}^{0}_{D}(\overline U)
      \longrightarrow
      \operatorname{Pic}^{0}(\overline U)
    \right],
\end{equation}
and realizes $H^1(U,\mathbb Z(1))$.  Huber--W\"ustholz prove that the de Rham--Betti realization of $1$-motives
is fully faithful and that its image is closed under subquotients
\cite[Theorem~9.14]{HuberWustholz}.  Closure under subquotients shows that the
specific Tate line in \eqref{eq:tate-twist} is the realization of a
$1$-motivic subobject; full faithfulness lifts its inclusion to the specific
morphism
\[
  \Q(0)\longrightarrow\operatorname{Pic}^{+}(U)_{\Q}
\]
whose realization is \eqref{eq:tate-twist}.

In the isogeny category this morphism is represented by a rational boundary
divisor $\Delta\in\operatorname{Div}^{0}_{D}(\overline U)\otimes\Q$ whose
class in $\operatorname{Pic}^{0}(\overline U)\otimes\Q$ is zero.  Choose an
integer $N>0$ which both clears denominators and kills the resulting torsion
Picard class.  Then the integral divisor $N\Delta$ has trivial associated line
bundle, so
\[
  N\Delta=\operatorname{div}(u)
\]
for a rational function whose zeros and poles lie in $D$; equivalently,
$u\in\Ocal(U)^\times$.  By \cref{prop:picard-convention}, the lattice generator represented by the
principal boundary divisor $\operatorname{div}(u)$ maps under de Rham
realization to the logarithmic class $[d\log u]$.  Compatibility of
the lifted morphism with de Rham realization therefore gives
\[
  N[\omega]=[d\log u]
  \quad\text{in }H^1_{\dR}(U/\overline{\Q}).
\]
Since $U$ is affine, algebraic de Rham cohomology in degree one is computed by
global forms \cite{HartshorneDR}.  Hence
\[
  N\omega-d\log u=dF
\]
for some $F\in\Ocal(U)$.  Therefore
\[
  L_0\simeq
  \Ecal^{-F/N}\otimes
  \left(\Ocal_U,d+\frac1N d\log u\right).
\]
The second factor has order dividing $N$, since its $N$th tensor power is
trivialized by $u^{-1}$.  Tensoring back by $A_0$ proves the theorem.
\end{proof}

\begin{remark}[Arithmetic scope]
Over $\C$, a topologically trivial rank-one algebraic connection can carry a
nonzero second-kind de Rham class and need not be an exponential times a finite
connection.  Huber--W\"ustholz excludes this non-Tate component over
$\overline{\Q}$.  For example, if $U=E\setminus\{0\}$ for an elliptic curve,
then $\operatorname{Div}^{0}_{\{0\}}(E)=0$, so
$\operatorname{Pic}^{+}(U)$ has no toric $\Q(0)$-subobject.  A nonzero
second-kind class over $\C$ therefore cannot arise from the Tate-line mechanism
used in the proof.
\end{remark}

\subsection{Generic saturation for the Kontsevich--Zagier category}

Let $\eta=\Spec k(S)$ be the generic point.  We write
\[
  \EMN(\eta):=\varinjlim_{U\subset S}\EMN(U),
\]
where objects and morphisms are identified after restriction to a smaller
nonempty smooth open.  Restriction is exact and tensor-compatible, so every
finite diagram in this filtered $2$-colimit has a common model and the generic
realization is an exact faithful tensor functor to algebraic differential
modules over $k(S)$.

Let $\Tcal_{\mathrm{KZ}}(\eta)$ denote the filtered union of all finitely
generated rigid Tannakian subcategories of $\EMN(\eta)$ generated by
dualizable restrictions of elementary relative Kontsevich--Zagier exponential
data.

\begin{theorem}[Generic rank-one saturation]\label{thm:global-RS}
For every $\lambda\in\overline{\Q}^{\times}$, every finitely generated rigid
Tannakian subcategory of $\Tcal_{\mathrm{KZ}}(\eta)$ is rank-one saturated.
\end{theorem}

\begin{proof}
Let $B$ be a tensor construction and subquotient in such a category, and let
$L$ be an invertible differential subquotient of $\rho_\lambda(B)$.  Choose a
common smooth open model on which $B$ and $L$ are lisse and choose a regular
geometric pre-Fourier representative $K_B$.  Such a representative exists
because the quotient functor defining $\EMN$ is essentially surjective on the
heart; changing it by a constant object does not change any nonzero
exponential realization.  Additive convolution and the six operations preserve
regular geometric origin.

After a further dense-open restriction, \cref{cor:rankone-finite} shows that
the Betti local system of $L$ has finite monodromy.  Applying
\cref{thm:arithmetic-kernel} on a dense affine open gives
\[
  L\simeq\Ecal^f\otimes A
\]
with $f$ algebraic and $A$ finite order.  At the fixed scale $\lambda$, the
graph object with potential $-f/\lambda$ realizes $\Ecal^f$.  The finite-order
character $A$ is cut out by its character idempotent in the finite \`etale
pushforward which trivializes it.  Their convolution product is an invertible
object in the generic exponential Nori category realizing $L$.  Hence $L$ has an ambient invertible lift.
\end{proof}

\begin{remark}[Models on dense opens]
Every finite packet is defined on a common dense lisse open.  If a rank-one
connection is defined on a larger smooth base, finite monodromy on a dense open
implies finite monodromy on the base because
$\pi_1(U^{\mathrm{an}})\twoheadrightarrow\pi_1(S^{\mathrm{an}})$.
\end{remark}

\section{Twisted fixed parts and Galois exactness}\label{sec:Tannaka}

Fix $\lambda\in\overline{\Q}^{\times}$ and let $\Tcal$ be a finitely generated
rigid Tannakian subcategory of $\Tcal_{\mathrm{KZ}}(\eta)$.  Choose a common
smooth connected open model $S$ on which every generator is $\lambda$-lisse;
all tensor constructions are then $\lambda$-lisse as well.  Fix a fiber functor
at a geometric point $s\in S(\C)$ and write
\[
  G_{\mot}=\Aut^\otimes(\omega_s|_{\Tcal})
\]
for the motivic Galois group.  Let $\Tcal_\nabla$ be the differential tensor
category generated by $\rho_\lambda(\Tcal)$ and write
\[
  G_{\diff}=\Aut^\otimes(\omega_s|_{\Tcal_\nabla}).
\]
The exact faithful tensor realization gives a closed immersion
$G_{\diff}\hookrightarrow G_{\mot}$.

\subsection{Semi-invariant spaces}

For a $G_{\mot}$-representation $V$ and a character
$\chi:G_{\diff}\to\mathbb G_m$, put
\[
  V^\chi
  =\{v\in V:h v=\chi(h)v\text{ for all }h\in G_{\diff}\}.
\]

\begin{proposition}[Motivic semi-invariants]\label{prop:semi-invariants}
Assume rank-one saturation.  For every $B\in\Tcal$ and every character $\chi$,
the subspace $\rho_\lambda(B)^\chi$ is the realization of a subobject of $B$ belonging
to $\Tcal$.
\end{proposition}

\begin{proof}
Let $L_\chi$ be the invertible differential object associated with $\chi$ and
choose an ambient invertible lift $\mathcal L_\chi$ by rank-one saturation.  Apply
\cref{thm:fixed-part} to the ambient object
$B\otimes\mathcal L_\chi^{-1}$ and let
\[
  \mathcal C_\chi
  :=(B\otimes\mathcal L_\chi^{-1})^{\mathrm{fix}}.
\]
Its realization is
\[
  (\rho_\lambda(B)\otimes L_\chi^{-1})^{G_{\diff}}.
\]
Tensor the inclusion $\mathcal C_\chi\hookrightarrow
B\otimes\mathcal L_\chi^{-1}$ by $\mathcal L_\chi$ and define
\[
  \mathcal W_\chi
  :=\im\!\left(
     \mathcal L_\chi\otimes\mathcal C_\chi\longrightarrow B
    \right).
\]
Exactness of realization gives
\[
  \rho_\lambda(\mathcal W_\chi)
  =L_\chi\otimes
    (\rho_\lambda(B)\otimes L_\chi^{-1})^{G_{\diff}}
  =\rho_\lambda(B)^\chi.
\]
Although $\mathcal L_\chi$ was chosen in the ambient category,
$\mathcal W_\chi$ is a subobject of $B$.  The Tannakian category generated by
$B$ is closed under subobjects, so $\mathcal W_\chi\in\Tcal$.
\end{proof}

\subsection{Normality}

We use Andre's criterion in the following form.

\begin{lemma}[Andre's normality criterion]\label{lem:Andre}
Let $H\subseteq G$ be affine algebraic groups over a field of characteristic zero.
If, for every finite-dimensional $G$-representation $V$ and every character
$\chi$ of $H$, the semi-invariant space $V^\chi$ is $G$-stable, then $H$ is
normal in $G$.
\end{lemma}

\begin{proof}
See \cite{AndreFixed} and the Tannakian formulation in \cite{Mostaed}.
\end{proof}

\begin{theorem}[Relative exponential Galois exactness]\label{thm:Galois-exact}
Assume rank-one saturation.  Let $\Tcal_0\subseteq\Tcal$ be the full
subcategory of constant objects and let $G_{\const}$ be its Tannakian group.
Then
\[
  1\longrightarrow G_{\diff}
   \longrightarrow G_{\mot}
   \longrightarrow G_{\const}
   \longrightarrow1
\]
is exact.
\end{theorem}

\begin{proof}
By \cref{prop:semi-invariants,lem:Andre}, $G_{\diff}$ is normal in
$G_{\mot}$.  By \cref{prop:constant-subquotients}, $\Tcal_0$ is a full tensor
subcategory stable under subquotients; hence its inclusion induces a faithfully
flat morphism $G_{\mot}\twoheadrightarrow G_{\const}$
\cite[Proposition~2.21]{DeligneMilne}.

An object of $\Tcal$ is trivial under $G_{\diff}$ if and only if its
$\lambda$-differential realization is trivial.  By
\cref{lem:constant-reflection-exp}, such an object is constant.  Conversely,
constant objects have trivial relative differential realization.  Therefore the
Tannakian quotient corresponding to $G_{\mot}/G_{\diff}$ is
$\Tcal_0$, proving exactness.
\end{proof}

\begin{corollary}[Generic split Fourier-geometric exactness]
If $\Tcal$ is finitely generated and split Fourier-geometric, then after
restriction to a suitable common dense open the exact sequence of
\cref{thm:Galois-exact} holds.
\end{corollary}

\section{The relative functional period theorem}\label{sec:periods}

\subsection{Comparison torsors and the constant-period base}

Let $K=k(S)$ and let $\omega_{\dR}$ and $\omega_{\Betti}$ denote the de Rham
and rapid-decay Betti fiber functors on $\Tcal$.  The relative motivic
comparison torsor is
\[
  \mathscr X_{\mot}
  =\Isom^\otimes(\omega_{\Betti}\otimes_k K,\omega_{\dR}),
\]
and its restriction to the constant subcategory is denoted $\mathscr X_0$.

Let $p_0$ be the actual constant comparison point and let
$Z_0\subseteq\mathscr X_0$ be its Zariski closure.  Put
\[
  A_0:=\Ocal(Z_0),
  \qquad
  R_0:=K\otimes_k A_0,
  \qquad
  F_0:=\operatorname{Frac}(A_0),
  \qquad
  K_0:=\operatorname{Frac}(R_0).
\]
The kernel of evaluation at $p_0$ is prime, so $A_0$ is a domain.  Since $S$
is geometrically integral over the algebraically closed field $k$, $R_0$ is a
domain as well.  We extend the derivations of $K$ to $K_0$ by declaring $F_0$
to be constant.

Define
\[
  \mathscr X_{\mathrm{rel}}
  :=\mathscr X_{\mot}\times_{\mathscr X_0}Z_0.
\]
The algebra $A_0$ is the coordinate ring of the Zariski closure of the actual
constant comparison point.

\begin{proposition}[Relative torsor]\label{prop:relative-torsor}
Under rank-one saturation, the morphism
\[
  \mathscr X_{\mathrm{rel}}\longrightarrow\Spec R_0
\]
is a $G_{\diff,R_0}$-torsor.  Consequently
$\Ocal(\mathscr X_{\mathrm{rel}})$ is faithfully flat and torsion-free over
$R_0$.
\end{proposition}

\begin{proof}
By \cref{thm:Galois-exact}, restriction of comparison isomorphisms to the
constant subcategory is the quotient of the $G_{\mot}$-torsor
$\mathscr X_{\mot}$ by the normal subgroup $G_{\diff}$.  Thus
$\mathscr X_{\mot}\to\mathscr X_0$ is a $G_{\diff}$-torsor.  Base change along
$Z_0\to\mathscr X_0$, followed by the scalar extension from $k$ to $K$, gives
the asserted torsor over $R_0$.  Torsors under affine group schemes are
faithfully flat over their base.
\end{proof}

\subsection{Comparison with Picard--Vessiot theory}

Let $\mathscr X_{\PV}$ be the Picard--Vessiot torsor of the differential tensor
category $\Tcal_\nabla$ over the differential field $K_0$, whose field of
constants contains $F_0$.  On a simply connected analytic open, the comparison
matrices of objects of $\Tcal$ generate the Picard--Vessiot algebra.  Evaluation
of motivic matrix coefficients at the same matrices gives a
$G_{\diff,K_0}$-equivariant closed immersion
\[
  \mathscr X_{\PV}
  \longrightarrow
  \mathscr X_{\mathrm{rel}}\times_{R_0}K_0.
\]

\begin{lemma}[Equivariant maps of torsors]\label{lem:torsor-map}
An equivariant morphism between two nonempty torsors under the same affine group
over a field is an isomorphism.
\end{lemma}

\begin{proof}
After a faithfully flat extension trivializing both torsors, the morphism becomes
an equivariant endomorphism of the group acting on itself by translation, hence a
translation and therefore an isomorphism.  Descent gives the result over the
original field.
\end{proof}

\begin{theorem}[Relative functional exponential-period theorem]
\label{thm:functional-period}
Assume rank-one saturation.  The generic fiber of $\mathscr X_{\mathrm{rel}}$
is the Picard--Vessiot torsor.  The analytic relative period map
\[
  \per^{\exp}_{S,\mathrm{rel}}:
  \Ocal(\mathscr X_{\mathrm{rel}})
  \longrightarrow\mathscr M(\widetilde S)
\]
is injective.
\end{theorem}

\begin{proof}
By \cref{prop:relative-torsor}, the generic fiber
$\mathscr X_{\mathrm{rel}}\times_{R_0}K_0$ and $\mathscr X_{\PV}$ are nonempty
torsors under the same group.  The preceding closed immersion is therefore an
isomorphism by \cref{lem:torsor-map}.

Let $x\in\Ocal(\mathscr X_{\mathrm{rel}})$ map to zero analytically.  After
tensoring with $K_0$, the torsor isomorphism identifies $x$ with an element of
the Picard--Vessiot algebra which is zero in the meromorphic solution field;
hence $x\otimes1=0$.  Thus $x$ is $R_0$-torsion.  Faithful flatness in
\cref{prop:relative-torsor} makes
$\Ocal(\mathscr X_{\mathrm{rel}})$ torsion-free over $R_0$, so $x=0$.
\end{proof}

\begin{corollary}[Functional relation interpretation]\label{cor:functional-rel}
Under rank-one saturation, every polynomial functional relation among the
relative exponential periods of $\Tcal$, with coefficients in the actual
constant exponential-period field, is induced by a tensor relation in $\Tcal$.
\end{corollary}

\begin{proof}
Polynomial expressions in matrix coefficients are matrix coefficients of finite
tensor constructions.  Injectivity of the comparison algebra therefore
identifies the kernel of functional evaluation with the tensor relations defining
the comparison torsor.
\end{proof}

\begin{remark}[Categorical and elementary period algebras]
\Cref{cor:functional-rel} gives completeness in the motivic formal-period algebra.
It does not by itself prove that this algebra has a minimal presentation solely by
linearity, elementary algebraic changes of variables, Fubini, and twisted Stokes.
The comparison between the categorical presentation and a chosen elementary
Kontsevich--Zagier presentation is a separate theorem.
\end{remark}

\section{Relation with Kontsevich--Zagier exponential periods}\label{sec:KZ}

Let $\Tcal_{\mathrm{KZ}}(\eta)$ be the generic relative exponential Nori
Tannakian envelope introduced before \cref{thm:global-RS}.  It is generated by
the restrictions of elementary data
\[
  (X\to S,Y,f,n,i),
  \qquad f:X\to\Aone\ \text{regular},
\]
and their tensor operations and duals.  Every finite packet has a common dense
lisse model.  The comparison theorems for exponential periods identify its
matrix coefficients with cohomological rapid-decay period functions
\cite{CHH,CommelinHuberII,FresanJossenExp}.

\begin{theorem}[Relative exponential Nori--Ayoub theorem]
\label{thm:exponential-NA}
For the canonical realization $\lambda=1$, the generic relative formal
exponential-period map of the filtered Kontsevich--Zagier exponential
Tannakian envelope is injective after the algebraic relations among constant
exponential periods have been imposed.  The same statement holds for every
$\lambda\in\overline{\Q}^{\times}$.  Equivalently, every algebraic functional
relation among a finite packet of relative exponential periods is induced, on
a common dense lisse open, by a tensor relation in the relative exponential
Nori category.
\end{theorem}

\begin{proof}
Every finite packet is contained in a finitely generated rigid Tannakian
subcategory of $\Tcal_{\mathrm{KZ}}(\eta)$.  By \cref{thm:global-RS}, that
subcategory is rank-one saturated at the chosen algebraic nonzero scale.  Choose
a common dense lisse model and apply \cref{thm:functional-period}.  Taking the filtered colimit over
finite packets proves the generic statement.  Since analytic functions agreeing
on a dense open agree wherever both sides are defined, this is the asserted
functional relation theorem.
\end{proof}

\subsection{Arithmetic specialization}

\Cref{thm:exponential-NA} does not solve the numerical exponential-period
conjecture.  It controls identities valid identically in the parameters.  A
relation created at an algebraic fiber requires an arithmetic lifting theorem.
In the $E$-function setting this role is played by Beukers' refinement of the
Siegel--Shidlovskii theorem; in period-valued, Gamma, and more general irregular
settings the corresponding specialization theorems remain independent problems.
The cyclotomic Gamma calculation in \cref{subsec:gamma-presentation} makes this
separation explicit.

\section{Applications and examples}\label{sec:examples}

\subsection{Polynomial exponential integrals}

For a finite packet of integrals
\[
  F_j(z)=\int_{\gamma_j}p_j(x)e^{-zq_j(x)}\,dx,
\]
the Fourier support is finite and algebraic.  Whenever the companion
polynomial-line analysis verifies that a finite cover splits the critical-value
support and that the regular coefficient systems have finite monodromy, the
generic split Fourier-geometric theorem applies.  The present result then gives
the categorical completeness of functional relations; specialization at an
algebraic value of $z$ remains the separate arithmetic $E$-function step.

\subsection{Cyclotomic Gamma packets and the Rohrlich--Lang boundary}
\label{subsec:gamma-presentation}

The main theorem gives categorical completeness for these rapid-decay periods.
For the standard one-parameter Gamma packet, the complete functional ideal can
also be computed directly.  The resulting Gamma-specific presentation is elementary and, to our knowledge,
has not been recorded in this form.

Fix an integer $N\geq2$ and distinct integers
\[
  1\leq m_1<\cdots<m_s<N.
\]
Put
\[
  r_j=\frac{m_j}{N},
  \qquad
  \gamma_j=\Gamma(r_j),
  \qquad
  C_\Gamma=\overline{\Q}(\gamma_1,\ldots,\gamma_s)\subset\C.
\]
For $\Re z>0$, define
\[
  G_j(z)
  =N\int_0^\infty x^{m_j-1}e^{-zx^N}\,dx,
\]
and continue the function to the universal cover $\widetilde{\Gm}$.
For a multi-index $\mathbf a=(a_1,\ldots,a_s)\in\mathbb N^s$, write
\[
  X^{\mathbf a}=\prod_{j=1}^sX_j^{a_j},
  \qquad
  \gamma^{\mathbf a}=\prod_{j=1}^s\gamma_j^{a_j},
  \qquad
  M(\mathbf a)=\sum_{j=1}^sa_jm_j.
\]

\begin{theorem}[Functional cyclotomic Gamma presentation]
\label{thm:functional-gamma-presentation}
One has
\[
  G_j(z)=\gamma_jz^{-m_j/N}.
\]
Let
\[
  \Phi:C_\Gamma(z)[X_1,\ldots,X_s]
  \longrightarrow \mathscr M(\widetilde{\Gm}),
  \qquad
  X_j\longmapsto G_j(z).
\]
Then
\begin{equation}\label{eq:gamma-kernel}
  \ker\Phi
  =\left(
      \gamma^{\mathbf b}X^{\mathbf a}
      -z^{-q}\gamma^{\mathbf a}X^{\mathbf b}
      \;:\;
      M(\mathbf a)-M(\mathbf b)=qN
    \right),
\end{equation}
where $\mathbf a,\mathbf b\in\mathbb N^s$ and $q\in\mathbb Z$.
Thus every polynomial functional relation in the packet is generated by the
Kummer-weight binomials in \eqref{eq:gamma-kernel}.
\end{theorem}

\begin{proof}
For positive real $z$, the substitution $u=zx^N$ gives
\[
  G_j(z)
  =z^{-m_j/N}\int_0^\infty u^{m_j/N-1}e^{-u}\,du
  =\gamma_jz^{-m_j/N}.
\]
Both sides are holomorphic on the right half-plane, so the identity holds there
and then on $\widetilde{\Gm}$ by analytic continuation.

Choose the global function
\[
  t=z^{-1/N}
\]
on $\widetilde{\Gm}$.  Then $t^N=z^{-1}$ and
\[
  G_j(z)=\gamma_jt^{m_j}.
\]
Since $\mu_N\subset\overline{\Q}\subset C_\Gamma$, the cyclic action
$t\mapsto\zeta t$ has fixed field $C_\Gamma(t^N)=C_\Gamma(z)$.  Hence
\[
  [C_\Gamma(t):C_\Gamma(z)]=N,
  \qquad
  1,t,\ldots,t^{N-1}
\]
is a basis over $C_\Gamma(z)$.

Let $I$ denote the ideal on the right-hand side of
\eqref{eq:gamma-kernel}.  Every displayed generator maps to zero, since
$M(\mathbf a)-M(\mathbf b)=qN$ implies
\[
  t^{M(\mathbf a)}=z^{-q}t^{M(\mathbf b)}.
\]
Thus $I\subseteq\ker\Phi$.

Conversely, let
\[
  P=\sum_{\mathbf a}c_{\mathbf a}(z)X^{\mathbf a}\in\ker\Phi.
\]
Write uniquely
\[
  M(\mathbf a)=q_{\mathbf a}N+\rho_{\mathbf a},
  \qquad
  0\leq\rho_{\mathbf a}<N.
\]
Then
\[
  0=\Phi(P)
   =\sum_{\rho=0}^{N-1}t^\rho
      \left(
        \sum_{\rho_{\mathbf a}=\rho}
        c_{\mathbf a}(z)\gamma^{\mathbf a}z^{-q_{\mathbf a}}
      \right).
\]
Linear independence of $1,t,\ldots,t^{N-1}$ gives, for every $\rho$,
\begin{equation}\label{eq:gamma-residue-coeff}
  \sum_{\rho_{\mathbf a}=\rho}
  c_{\mathbf a}(z)\gamma^{\mathbf a}z^{-q_{\mathbf a}}=0.
\end{equation}
Fix one exponent vector $\mathbf a_0$ in a nonempty residue class.  Modulo $I$,
\[
  X^{\mathbf a}
  \equiv
  z^{-(q_{\mathbf a}-q_{\mathbf a_0})}
  \frac{\gamma^{\mathbf a}}{\gamma^{\mathbf a_0}}
  X^{\mathbf a_0}
  \qquad
  (\rho_{\mathbf a}=\rho_{\mathbf a_0}).
\]
Therefore the part of $P$ in this residue class is congruent modulo $I$ to
\[
  \frac{z^{q_{\mathbf a_0}}X^{\mathbf a_0}}
       {\gamma^{\mathbf a_0}}
  \left(
    \sum_{\rho_{\mathbf a}=\rho_{\mathbf a_0}}
    c_{\mathbf a}(z)\gamma^{\mathbf a}z^{-q_{\mathbf a}}
  \right),
\]
which vanishes by \eqref{eq:gamma-residue-coeff}.  Summing over the residue
classes gives $P\in I$.
\end{proof}

For the complete packet, the general binomial ideal has a particularly small
set of generators.

\begin{corollary}[Complete fixed-denominator packet]
\label{cor:complete-gamma-packet}
For
\[
  G_m(z)=\Gamma(m/N)z^{-m/N},
  \qquad 1\leq m<N,
\]
put $\gamma_m=\Gamma(m/N)$.  The kernel of
\[
  C_{\Gamma,N}(z)[X_1,\ldots,X_{N-1}]
  \longrightarrow\mathscr M(\widetilde{\Gm}),
  \qquad X_m\longmapsto G_m(z),
\]
where
\[
  C_{\Gamma,N}=\overline{\Q}(\gamma_1,\ldots,\gamma_{N-1}),
\]
is generated by
\[
  \gamma_1^mX_m-\gamma_mX_1^m
  \qquad(2\leq m<N)
\]
and
\[
  X_1^N-\gamma_1^Nz^{-1}.
\]
\end{corollary}

\begin{proof}
The first family eliminates every $X_m$ in favor of $X_1$.  The resulting
quotient is
\[
  C_{\Gamma,N}(z)[X_1]/(X_1^N-\gamma_1^Nz^{-1}).
\]
The map $X_1\mapsto\gamma_1t$ embeds this quotient into
$C_{\Gamma,N}(t)$ because $1,t,\ldots,t^{N-1}$ are linearly independent
over $C_{\Gamma,N}(z)$.  Hence the
displayed relations generate the full kernel.
\end{proof}

For example, when $N=3$,
\[
  G_1(z)=\Gamma(1/3)z^{-1/3},
  \qquad
  G_2(z)=\Gamma(2/3)z^{-2/3},
\]
and the complete functional ideal is generated by
\[
  \Gamma(1/3)^2G_2(z)-\Gamma(2/3)G_1(z)^2
\]
and
\[
  G_1(z)^3-\Gamma(1/3)^3z^{-1}.
\]

\begin{remark}[Relation with the Rohrlich--Lang conjectures]
\label{rem:gamma-rohrlich-lang}
Rohrlich's conjecture predicts that the multiplicative algebraic relations among
special rational Gamma values and $2\pi i$ are generated by the standard
translation, reflection, and multiplication identities.  Lang's strengthening
makes the corresponding prediction for all polynomial algebraic relations
\cite[\S1.4]{AndersonBrownawellPapanikolas}; see also \cite{RivoalGamma}.
Koblitz and Ogus proved the algebraicity implication in the
Deligne--Koblitz--Ogus criterion for Gamma monomials; the conjectural converse
is a form of Rohrlich's conjecture \cite{KoblitzOgus}.

\Cref{thm:functional-gamma-presentation} is a relative functional statement,
not a result on those constant-value conjectures.  Its coefficient field
$C_\Gamma$ is the actual field generated by the chosen Gamma constants, so all
relations among those constants have already been imposed.  Before extending
scalars from $\overline{\Q}$ to $C_\Gamma$, the same residue-class argument
reduces the unknown part of the functional kernel to linear relations among the
Gamma monomials $\gamma^{\mathbf a}$.  Thus the functional variable contributes
only the Kummer-weight binomials above; the remaining arithmetic problem is
the relation problem among the chosen Gamma constants, together with $2\pi i$
in the usual global formulation.  In this sense,
the theorem gives a fixed-denominator functional analogue of the
Rohrlich--Lang picture while leaving its classical numerical core open.
\end{remark}

\subsection{\texorpdfstring{Nilsson--$E$ packets}{Nilsson--E packets}}

A finite packet of functions
\[
  z^\alpha(\log z)^mF(z),
  \qquad \alpha\in\Q,
\]
with $F$ an $E$-function has finite Kummer and unipotent monodromy at the regular
singularity.  Monodromy projectors and finite differences provide a concrete
split model for separating the Kummer and logarithmic components.  Once the
packet is placed in the relative exponential Nori envelope, the generic
functional theorem applies.  Relations at a fixed algebraic value with
period-valued coefficients remain arithmetic and are not consequences of the
functional theorem.

\section{Concluding perspective}

The relative exponential Nori category admits the same functional-period
structure as the ordinary motivic setting.  The fixed-part theorem identifies
constant exponential submotives, while the Fourier nearby-cycle and
$1$-motivic arguments control differential characters.  Together they give the
relative Galois exact sequence and identify the generic comparison torsor with
the Picard--Vessiot torsor.

For the filtered generic Kontsevich--Zagier exponential envelope, the formal
functional period map is therefore injective at every algebraic nonzero
Fourier scale.  The cyclotomic Gamma calculation gives an explicit instance:
over the field generated by the constant Gamma values, the complete functional
ideal is the Kummer binomial ideal of
\cref{thm:functional-gamma-presentation}, while the classical
Rohrlich--Lang problem remains entirely in the constants.  Arithmetic relations
created by specialization and an elementary presentation of the general motivic
relation algebra are separate problems.

\section*{AI-assisted research disclosure}
The authors used AI-assisted tools for literature retrieval, drafting, and proof
checking, and take responsibility for the manuscript.  The AI system is not an
author.

\begin{thebibliography}{99}

\bibitem{AndersonBrownawellPapanikolas}
G.~W. Anderson, W.~D. Brownawell, and M.~A. Papanikolas,
\emph{Determination of the algebraic relations among special $\Gamma$-values in positive characteristic},
Ann. of Math. (2) \textbf{160} (2004), no.~1, 237--313.

\bibitem{AndreFixed}
Y.~Andre,
\emph{Mumford--Tate groups of mixed Hodge structures and the theorem of the fixed part},
Compositio Math. \textbf{82} (1992), no.~1, 1--24.

\bibitem{AyoubPeriods}
J.~Ayoub,
\emph{Periods and the conjectures of Grothendieck and Kontsevich--Zagier},
Eur. Math. Soc. Newsl. (2014), no.~91, 12--18.

\bibitem{AyoubRelative}
J.~Ayoub,
\emph{Une version relative de la conjecture des periodes de Kontsevich--Zagier},
Ann. of Math. (2) \textbf{181} (2015), no.~3, 905--992.

\bibitem{AyoubBarbieriViale}
J.~Ayoub and L.~Barbieri-Viale,
\emph{Nori $1$-motives},
Math. Ann. \textbf{361} (2015), no.~1--2, 367--402.

\bibitem{BarbieriVialeSrinivas}
L.~Barbieri-Viale and V.~Srinivas,
\emph{Albanese and Picard $1$-motives},
M\'em. Soc. Math. Fr. (N.S.) \textbf{87} (2001), vi+104 pp.

\bibitem{BBD}
A.~A. Beilinson, J.~Bernstein, and P.~Deligne,
\emph{Faisceaux pervers},
Asterisque \textbf{100} (1982).

\bibitem{BakkerTsimerman}
B.~Bakker and J.~Tsimerman,
\emph{Functional transcendence of periods and the geometric Andre--Grothendieck period conjecture},
Forum Math. Sigma \textbf{13} (2025), e97.

\bibitem{CHH}
J.~Commelin, P.~Habegger, and A.~Huber,
\emph{Exponential periods and o-minimality},
arXiv:2007.08280, version 4, 2026.

\bibitem{CommelinHuberII}
J.~Commelin and A.~Huber,
\emph{Exponential periods and o-minimality II},
arXiv:2007.08290.

\bibitem{FresanJossenExp}
J.~Fresan and P.~Jossen,
\emph{Exponential motives},
preliminary manuscript, version consulted August 2026, available from the authors.

\bibitem{GallauerPepin}
M.~Gallauer and S.~Pepin Lehalleur,
\emph{Exponentiation of coefficient systems and exponential motives},
arXiv:2211.17247.

\bibitem{HuberWustholz}
A.~Huber and G.~W\"ustholz,
\emph{Transcendence and Linear Relations of $1$-Periods},
Cambridge Tracts in Mathematics, Cambridge University Press, 2022.

\bibitem{Hien}
M.~Hien,
\emph{Periods for flat algebraic connections},
Invent. Math. \textbf{178} (2009), no.~1, 1--22.

\bibitem{HienRoucairol}
M.~Hien and C.~Roucairol,
\emph{Integral representations for solutions of exponential Gauss--Manin systems},
Bull. Soc. Math. France \textbf{136} (2008), no.~4, 505--532.

\bibitem{IvorraMorel}
F.~Ivorra and S.~Morel,
\emph{The four operations on perverse motives},
J. Eur. Math. Soc. \textbf{26} (2024), 4191--4272.


\bibitem{DeligneHodgeIII}
P.~Deligne,
\emph{Th\'eorie de Hodge. III},
Publ. Math. Inst. Hautes \'Etudes Sci. \textbf{44} (1974), 5--77.

\bibitem{DeligneMilne}
P.~Deligne and J.~S. Milne,
\emph{Tannakian categories},
in \emph{Hodge Cycles, Motives, and Shimura Varieties},
Lecture Notes in Math. \textbf{900}, Springer, 1982, pp.~101--228.

\bibitem{DouaiSabbah}
A.~Douai and C.~Sabbah,
\emph{Gauss--Manin systems, Brieskorn lattices and Frobenius structures (I)},
Ann. Inst. Fourier (Grenoble) \textbf{53} (2003), no.~4, 1055--1116.

\bibitem{JacobsenComparison}
E.~Jacobsen and L.~Terenzi,
\emph{A comparison of categories of Nori motivic sheaves},
arXiv:2509.21476v1, 2025.

\bibitem{KatzESDE}
N.~M.~Katz,
\emph{Exponential Sums and Differential Equations},
Annals of Mathematics Studies \textbf{124}, Princeton University Press, 1990.

\bibitem{KontsevichZagier}
M.~Kontsevich and D.~Zagier,
\emph{Periods},
in \emph{Mathematics Unlimited--2001 and Beyond}, Springer, 2001, pp.~771--808.

\bibitem{KoblitzOgus}
N.~Koblitz and A.~Ogus,
\emph{Algebraicity of some products of values of the $\Gamma$ function},
appendix to P.~Deligne, \emph{Valeurs de fonctions $L$ et p\'eriodes d'int\'egrales},
Proc. Sympos. Pure Math. \textbf{33}, part~2, Amer. Math. Soc., 1979,
pp.~343--346.

\bibitem{LamLitt}
Y.-H.~J. Lam and D.~Litt,
\emph{Geometric local systems on the projective line minus four points},
Compos. Math. \textbf{161} (2025), no.~3, 536--554.

\bibitem{Mostaed}
A.~Mostaed,
\emph{On the invariant cycle theorem for families of Nori motives},
arXiv:2206.05582.

\bibitem{HartshorneDR}
R.~Hartshorne,
\emph{On the de Rham cohomology of algebraic varieties},
Publ. Math. Inst. Hautes Etudes Sci. \textbf{45} (1975), 5--99.

\bibitem{HTT}
R.~Hotta, K.~Takeuchi, and T.~Tanisaki,
\emph{$D$-Modules, Perverse Sheaves, and Representation Theory},
Progress in Mathematics \textbf{236}, Birkh\"auser, 2008.

\bibitem{RivoalGamma}
T.~Rivoal,
\emph{On the arithmetic nature of the values of the Gamma function, Euler's constant, and Gompertz's constant},
Michigan Math. J. \textbf{61} (2012), no.~2, 239--254.

\bibitem{SaitoMHM}
M.~Saito,
\emph{Mixed Hodge modules},
Publ. Res. Inst. Math. Sci. \textbf{26} (1990), no.~2, 221--333.

\bibitem{SabbahTwistor}
C.~Sabbah,
\emph{Polarizable Twistor $D$-Modules},
Ast\'erisque \textbf{300}, Soci\'et\'e Math\'ematique de France, 2005.

\bibitem{SabbahMonodromy}
C.~Sabbah,
\emph{Monodromy at infinity and Fourier transform II},
Publ. Res. Inst. Math. Sci. \textbf{42} (2006), no.~3, 803--835.

\bibitem{SabbahParameters}
C.~Sabbah,
\emph{A short proof of a theorem of Cotti, Dubrovin and Guzzetti},
Port. Math. \textbf{79} (2022), no.~1, 29--43;
arXiv:2103.16878v3.

\bibitem{Snodgrass}
B.~Snodgrass,
\emph{Periods of $E$-operators},
arXiv:2608.06005, 2026.

\bibitem{Terenzi}
L.~Terenzi,
\emph{Tensor structure on perverse Nori motives},
Ann. K-Theory \textbf{11} (2026), no.~1, 47--170; arXiv:2401.13547v3.

\bibitem{Tubach}
S.~Tubach,
\emph{On the Nori and Hodge realisations of Voevodsky motives},
Compos. Math. \textbf{161} (2025), no.~9, 2155--2201; arXiv:2309.11999v4.

\bibitem{Virk}
R.~Virk,
\emph{Realizations of exponential sheaves and Fourier transform},
arXiv:2605.14904v4, 2026.

\end{thebibliography}

\end{document}
